\section{Bijective transforms}
\label{app:bijective}

The transform $g_\theta$ must be invertible so that we can always recover the adapted representation needed for decoding (\secref{kron-method}). We guarantee this invertibility by structuring the channel transformation matrix $\matrix{C}$. Specifically, we write $\matrix{C}$ as
\begin{align}
   \label{eq:c-factor}
   \matrix{C} = \matrix{R} \, \mathrm{diag}(\sigma_1, \ldots, \sigma_{n_d}),
   \qquad \sigma_1, \ldots, \sigma_{n_d} > 0,
\end{align}
where $\matrix{R}$ is a rotation matrix (orthogonal) and $\mathrm{diag}(\sigma_1,\ldots,\sigma_{n_d})$ is a diagonal matrix with all positive entries. Both of these components are invertible, i.e., a diagonal matrix with positive values is inverted simply by taking the reciprocal of each entry, and a rotation is inverted by transposing it. Therefore, their product $\matrix{C}$ is also invertible.

To make sure that $\matrix{R}$ stays a true rotation during training, we construct it by exponentiating a skew-symmetric matrix:
\begin{align}
   \label{eq:q-expm}
   \matrix{R} = \exp(\matrix{P} - \matrix{P}^{\top})
\end{align}
where $\matrix{P}$ is any square matrix of size $n_d \times n_d$ (with no constraints). The difference $\matrix{P} - \matrix{P}^{\top}$ is always skew-symmetric by construction, and exponentiating any skew-symmetric matrix produces an orthogonal (rotation) matrix~\citep{lezcano2019cheap}, which lets us freely optimize $\matrix{P}$ as unconstrained parameters.

The full inverse of $g_\theta$ is given by
\begin{align}
   \label{eq:transform-inverse}
   g_\theta^{-1}(\matrix{V}') = a_\ell^{-1} \matrix{V}' \matrix{C}^{-1},
\end{align}
where $a_\ell > 0$ is a scalar learned separately for each layer (\secref{kron-method}). Since $g_\theta$ is linear, operates along the channel axis, and uses a simple per-layer scaling, it commutes with mean-pooling over the rank axis.
In other words, it produces the same result if we pool the rank dimension before or after applying $g_\theta$. Therefore, we can use $g_\theta$, trained on full-rank tensors, on mean-pooled vectors as well.

We initialize $g_\theta$ to the identity (i.e., no transformation initially) and keep all parameters of the base language model frozen. Training for $g_\theta$ uses $42{,}332$ examples from OpenAlex (see Appendix~\ref{app:data}) paired up as anchor, positive, easy negative, and hard negative.
The easy and hard negatives follow the sampling protocol of \texttt{SPECTER}. An easy negative is a paper drawn uniformly at random from the corpus, excluding the anchor and the papers it cites. A hard negative is a paper cited by the positive but not by the anchor, i.e., at distance two from the anchor. The optimization uses an InfoNCE contrastive loss with temperature $0.05$, in which each anchor is scored against its own hard and easy negative together with the other $255$ positives in the batch. We train for $8{,}000$ Adam steps with batch size $256$ at a learning rate of $10^{-3}$ and no added regularization.